\documentclass[12pt]{article}
\usepackage{indentfirst}
\usepackage{amsmath, amssymb, amsthm}
\usepackage[margin=2.5cm]{geometry}

\title{Understanding Linear Algebra}
\author{Shengran Sun}
\date{\today}

\begin{document}
\maketitle

\section{Why Matrix Multiplication is defined as that form?}

Let's consider 
\begin{align*}
S &= aU + bV \\
T &= cU + dV \\
U &= eX + fY \\
V &= gX + hY
\end{align*}
\indent Now we determine $S$ snd $T$ in terms of $X$ and $Y$.\\
\indent We get
\begin{align*}
S &= (ae + bg)X + (af + bh)Y \\
T &= (ce + dg)X + (cf + dh)Y
\end{align*}
\indent Here we might have noticed that \\
\indent if we write the equations in matrix form:
\[
\begin{pmatrix} S \\ T \end{pmatrix} = 
\begin{pmatrix} a & b \\ c & d \end{pmatrix}
\begin{pmatrix} U \\ V \end{pmatrix}
\]
\[
\begin{pmatrix} U \\ V \end{pmatrix} = 
\begin{pmatrix} e & f \\ g & h \end{pmatrix}
\begin{pmatrix} X \\ Y \end{pmatrix}
\]
\indent so that we can get \\
\[
\begin{pmatrix} S \\ T \end{pmatrix} = 
\begin{pmatrix} a & b \\ c & d \end{pmatrix}
\begin{pmatrix} e & f \\ g & h \end{pmatrix}
\begin{pmatrix} X \\ Y \end{pmatrix}
\]
\indent Hence we define the matrix multiplication as \\
\[
\begin{pmatrix} a & b \\ c & d \end{pmatrix}
\begin{pmatrix} e & f \\ g & h \end{pmatrix} =
\begin{pmatrix} ae+bg & af+bh \\ ce+dg & cf+dh \end{pmatrix}
\]

\section{Why Determinant is defined as that form?}
Consider a unit square with coordinates
\[
\begin{pmatrix}
    0 & 1 & 1 & 0 \\
    0 & 0 & 1 & 1
\end{pmatrix}
\]
\indent which places the coordinates of the four vertices anticlockwisely.\\
\indent Now we perform a transformation $A$ on it 
\[
\begin{pmatrix} a & b \\ c & d \end{pmatrix}
\begin{pmatrix} 0 & 1 & 1 & 0 \\ 0 & 0 & 1 & 1 \end{pmatrix}
\]
\indent The result is
\[
\begin{pmatrix} 0 & a & a+b & b \\ 0 & c & c+d & d \end{pmatrix}
\]
\indent which also places the coordinates of the four vertices anticlockwisely.\\
\indent We can find the area of the new shape is 
\[
ad - bc
\]
\indent which is then called the determinant of matrix $A$. \\
\indent However, if a matrix has negative determinant, then the way that it places the coordinates of the four vertices will be clockwise. \\
\indent For example, \\
\[
B = \begin{pmatrix} b & d \\ a & c \end{pmatrix}
\]
\indent has determinant $-(ad-bc)$. \\
\indent If the determinant of some transformation is $0$, then the 'dimension' of the original geometric object is 'reduced'. \\
\indent Actually, the determinant is zero (singular) iff the rows/columns are linearly dependent because
\begin{align*}
\det(A) &= ad - bc = 0 \\
\Rightarrow \indent ad &= bc \\
\Rightarrow \indent \frac{a}{b} &= \frac{c}{d} \indent \text{here we just assume a, b, c, d are not zero}\\
\end{align*}

\section{Inverse}
\subsection{Claim1}
The left inverse $L$ and right inverse $R$ of a matrix $A$ are equal. \\
\indent Proof: \\
\begin{align*}
(LA)R &= IR = R \\
&= L(AR) \indent \text{by associativity of matrix multiplication} \\
&= LI = L \\
\Rightarrow L = R
\end{align*}

\subsection{Claim2}
A square matrix $A$ has a left inverse $\Leftrightarrow$ it has a right inverse. \\
\indent Proof: \\
\begin{align*}
(\Rightarrow) \indent & LA = I \\
&\det(LA) = \det(L)\det(A) = \det(I) = 1 \\
& \Rightarrow \det(A) \ne 0 \Rightarrow A \text{ is invertible} \Rightarrow A^{-1} \text{ exists}. \\
& \Rightarrow (LA)A^{-1} = IA^{-1} \\
& \text{Also, } (LA)A^{-1} = L(AA^{-1}) = LI \\
& \Rightarrow IA^{-1} = LI \Rightarrow L = A^{-1} \\
(\Leftarrow) \indent & AR = I \\
& \det(AR) = \det(A)\det(R) = \det(I) = 1 \\
& \Rightarrow \det(A) \ne 0 \Rightarrow A \text{ is invertible} \Rightarrow A^{-1} \text{ exists}. \\
& \Rightarrow A^{-1}(AR) = A^{-1}I \\
& \text{Also, } A^{-1}(AR) = (A^{-1}A)R = IR \\
& \Rightarrow A^{-1}I = IR \Rightarrow R = A^{-1}\\
\end{align*}

\section{Why the inverse of a 2*2 matrix is defined as that form?}
Given a 2*2 matrix $M$, we can determine its inverse $M^{-1}$ by using $MM^{-1} = I$. \\
After that, we can find it corresponds to the defined form. \\
\begin{align*}
\begin{pmatrix} a & b \\ c & d \end{pmatrix} ^{-1} = \frac{1}{\det(M)} \begin{pmatrix} d & -b \\ -c & a \end{pmatrix}
\end{align*}
Reversing a transformation is exactly convert basis from the original to the priginal vectors before being transformed.



\end{document}